\subsection{为什么 2025 的 Agent 浪潮必须同时谈 Safety 与 Security}

slides 1--8 先给出课程定位。Frontier AI 能力快速提升，Agent 又把能力连接到真实工具与业务；风险既包括误用、诈骗、错误信息、隐私伤害和网络攻击，也包括系统自身在错误目标下造成的外部伤害。Safety 关注系统可能伤害外部环境，Security 关注攻击者如何伤害系统或借系统伤害他人，两者在 Agent 场景中高度交叉。

\lecturefigure{slides-images/slide-001.png}{官方 slide 1：Towards Building Safe and Secure Agentic AI}
\lecturefigure{slides-images/slide-002.png}{官方 slide 2：Frontier AI 的快速进展}
\lecturefigure{slides-images/slide-004.png}{官方 slide 4：2025 is the year of Agents}
\lecturefigure{slides-images/slide-005.png}{官方 slide 5：AI 风险的宽谱分类}
\lecturefigure{slides-images/slide-006.png}{官方 slide 6：有价值技术总会吸引攻击者}
\lecturefigure{slides-images/slide-007.png}{官方 slide 7：AI Safety 与 AI Security 的区别和交集}
\lecturefigure{slides-images/slide-008.png}{官方 slide 8：安全创新的总体目标}

\begin{knowledgebox}{读图：Safety 与 Security 不能分成两个团队各自收尾}
若模型因目标错配主动执行危险动作，这是 safety 问题；若攻击者通过网页内容诱导同一动作，这是 security 问题；两者在后果层可能完全相同。资产、权限、行动和恢复必须共享一套风险语言，否则安全团队只看越狱，产品团队只看误操作，攻击链会从组织缝隙穿过。
\end{knowledgebox}

\begin{warningbox}{课堂提示：默认存在主动攻击者}
历史经验表明，能力越有用，攻击收益越高。Agent 安全不能只用“正常用户平均表现”推断风险，而要假设攻击者会持续搜索 prompt、memory、tool schema 和业务流程中的组合漏洞，并会根据防御响应调整策略。
\end{warningbox}

